\chapter{Source digests I: foundations and Lyubich's geometry}
\label{ch:digests-foundations}

This chapter and the next state, source by source, the definitions and the
theorems on which the approaches to MLC rest.  Each section gives the
definitions the source uses, the results it states, the status the source
assigns to each result, and the other sources it depends on.  The statements
are those of the source, up to notation.  Where a statement is proved in this
review rather than in the source, this is said explicitly.  Lean
correspondences, when they exist, are given at the end of a section; a Lean
declaration is never a substitute for the mathematical statement.

\section{Notation and the quadratic family}

For $c\in\C$ let $f_c(z)=z^2+c$, and write $f_c^{n}$ for its $n$-th iterate.
Put
\begin{align*}
 \Kc&=\{z\in\C:\ \sup_{n\ge0}|f_c^n(z)|<\infty\},\\
 \Jc&=\partial\Kc,\\
 \M&=\{c\in\C:\ 0\text{ has a bounded forward orbit under }f_c\}.
\end{align*}
The classical equivalence $c\in\M\Leftrightarrow\Kc$ is connected is the
starting point of the subject; the statements below use both descriptions.

\begin{definition}[Quadratic family, critical orbit, MLC]
\label{def:mlc-statement}
The critical orbit of $f_c$ is $p_n(c)=f_c^n(0)$.  A point $c\in\M$ is a point of
local connectivity of $\M$ if $\M$ is locally connected at $c$.  The Mandelbrot
set is locally connected, $\M$ being $\LocallyConnectedSpace$, if this holds at
every point.  This last statement is the MLC conjecture.
\end{definition}

\begin{theorem}[Douady--Hubbard; compactness and connectedness]
\label{thm:dh-compact-connected}
$\M$ is a compact connected subset of $\C$.  Moreover $\M\subset\{|c|\le2\}$.
\end{theorem}

\begin{remark}
In the review, the inclusion $\M\subset\{|c|\le2\}$ is proved in
Chapter~\ref{ch:quadratic-foundations}, and the compactness of finite outer
approximations is formalized in Lean.  Connectedness of $\M$ is cited from
Douady and Hubbard~\cite{DH}; the review does not formalize it.
\end{remark}

\begin{theorem}[Real slice; proved in this review]
\label{thm:real-slice}
For $x\in\R$ one has $x\in\M$ if and only if $-2\le x\le\tfrac14$.
\end{theorem}

\begin{proof}
Let $r_n=f_x^n(0)$ for real $x$; then $r_{n+1}=r_n^2+x$ is real.

(i) \emph{Escape for $|x|>2$.}  Let $a=|x|>2$ and $\delta=a(a-2)>0$.  Then
$|f_x(z)|\ge|z|^2-a$ always, and $t^2-a-t\ge a^2-2a=\delta$ whenever $t\ge a$,
since $t^2-t-a^2+a=(t-a)(t+a-1)\ge0$.  Induction gives
$|f_x^{n+1}(0)|\ge a+n\delta$, so the orbit is unbounded.  The same argument
works for every complex $c$ with $|c|>2$; this is the escape part of
Chapter~\ref{ch:quadratic-foundations}.

(ii) \emph{Growth for $x>\tfrac14$.}  One has $r_{n+1}-r_n-(x-\tfrac14)=(r_n-\tfrac12)^2\ge0$,
so $r_n\ge n(x-\tfrac14)\to\infty$.

(iii) \emph{Invariant interval for $-2\le x\le\tfrac14$.}  Let $s=\sqrt{1-4x}\in[0,3]$
and $\beta=(1+s)/2$.  Then $\beta^2+x=\beta$, and $x+\beta=(3-s)(1+s)/4\ge0$, so
$-\beta\le x$.  If $|t|\le\beta$ then $t^2+x\le\beta^2+x=\beta$ and
$t^2+x\ge x\ge-\beta$.  Thus $[-\beta,\beta]$ is invariant, contains $0$, and
$\beta\le2$.  The orbit of $0$ is bounded by $2$.

Parts (i)--(iii) give the equivalence.
\end{proof}

\begin{remark}
Theorem~\ref{thm:real-slice} is classical; it is proved here by the argument
above, and the proof is formalized in
\leanref{Mlc/Quadratic/RealSlice.lean}{ofReal_mem_mandelbrotSet_iff}.
The Lean module also contains the linear escape estimate
\leanref{Mlc/Quadratic/RealSlice.lean}{norm_orbit_ge_of_two_lt_norm}.
\end{remark}

\section{Results on local connectivity cited in the sources}

The sources use the following historical results as inputs.  They are cited,
not proved here.  The statements follow the account in Lyubich's
paper~\cite{LyubichGeometry}, whose abstract records the attribution.

\begin{theorem}[Douady--Hubbard; quasi-hyperbolic points]
\label{thm:dh-quasihyperbolic}
$\M$ is locally connected at every quasi-hyperbolic parameter in the sense of
Douady and Hubbard~\cite{DH}.  Status: cited.
\end{theorem}

\begin{remark}
The definition of quasi-hyperbolic parameters is not restated here; the
review's definition list does not yet include it.  This is one of the gaps listed
in Section~\ref{sec:digests-gaps}.
\end{remark}

\begin{theorem}[Yoccoz; boundaries of hyperbolic components]
\label{thm:yoccoz-hyperbolic-boundary}
$\M$ is locally connected at every point of the boundary of a hyperbolic
component.  Status: cited, attributed to Yoccoz in~\cite{LyubichGeometry}.
\end{theorem}

\begin{theorem}[Yoccoz; finitely renormalizable parameters]
\label{thm:yoccoz-finite-renormalization}
MLC holds at every finitely renormalizable parameter $c\in\M$.  Status: cited,
attributed to Yoccoz in~\cite{LyubichGeometry} and in the survey~\cite{DudkoMLC}.
The renormalization used here is the Douady--Hubbard renormalization of
polynomials; its definition belongs to Chapter~\ref{ch:quadratic-foundations} and to the
survey~\cite{DudkoMLC}.
\end{theorem}

\begin{remark}
Theorems~\ref{thm:dh-quasihyperbolic}, \ref{thm:yoccoz-hyperbolic-boundary} and
\ref{thm:yoccoz-finite-renormalization} are the reason that MLC is open only at
infinitely renormalizable parameters outside the cases treated in the later
sources.  The survey~\cite{DudkoMLC} gives the same reduction in its Introduction.
\end{remark}

\section{Lyubich 1993: moduli, rigidity, and local connectivity}
\label{sec:digest-lyubich-1993}

Lyubich's paper~\cite{LyubichGeometry} is the first systematic use of complex
a priori bounds for a class of infinitely renormalizable quadratic polynomials.
Its stated aim is MLC at some infinitely renormalizable parameters, under two
combinatorial assumptions (C1, C2 below).

\subsection*{Definitions}

\begin{definition}[Quadratic-like map]
\label{def:quadratic-like}
A quadratic-like map is a proper holomorphic map $f\colon U'\to U$ of degree two
between topological disks with $\overline{U'}\subset U$.  Its filled Julia set is
$K_f=\{z\in U':\ f^n(z)\in U'\ \forall n\}$, and $J_f=\partial K_f$.
\end{definition}

\begin{definition}[Principal nest and quadratic-like renormalization]
\label{def:digest-principal-nest}
The principal nest is a sequence $V^0\supset V^1\supset\cdots$ of critical
puzzle pieces, where $V^{n+1}$ is obtained by pulling $V^n$ back along the first
return of the critical point.  If the first return map $f^{\pi(t)}\colon V^{t+1}\to V^t$
is itself quadratic-like, then $f$ is quadratic-like renormalizable and
$Rf=f^{\pi(t)}\colon V^{t+1}\to V^t$.  Iterating gives $R^nf$.
\end{definition}

\begin{definition}[Central returns and combinatorial type]
\label{def:combinatorial-type}
A return at level $n$ is central if the first return map sends the critical point
into $V^n$.  Consecutive central returns form a cascade.  The combinatorial type
of $f$ is the number $\kappa(f)$ of levels with non-central return.
\end{definition}

\begin{definition}[Combinatorial rotation number]
\label{def:rotation-number-lyubich}
If $p$ external rays land at the dividing fixed point $\alpha$ and are cyclically
permuted with rotation number $q/p$, the combinatorial rotation number of $\alpha$
is $q/p$.  Assumption (C1) below concerns these numbers.
\end{definition}

\begin{definition}[Definite modulus]
\label{def:definite-modulus}
A quadratic-like map has definite modulus $\bar\mu>0$ if
$\modulus(U\setminus\overline{U'})\ge\bar\mu$.
\end{definition}

\begin{definition}[Class $\mathcal F$]
\label{def:class-F}
$\mathcal F$ consists of the quadratic-like maps with infinitely many
quadratic-like renormalizations such that
\begin{itemize}
\item[(C1)] all renormalization hybrid classes lie in finitely many truncated secondary
  Mandelbrot limbs, so that the combinatorial rotation numbers are bounded;
\item[(C2)] every $\kappa(R^mf)$ is sufficiently high, depending on those limbs.
\end{itemize}
The threshold in (C2) is not quantified in the source.
\end{definition}

\begin{definition}[Hybrid equivalence and combinatorial equivalence]
\label{def:hybrid-equivalence}
Two quadratic-like maps are hybrid equivalent if they are conjugate by a
quasiconformal map $h$ with $\bar\partial h=0$ almost everywhere on the filled
Julia set.  Two maps are combinatorially equivalent if they induce the same
equivalence relation on rational external-ray arguments.  A quadratic polynomial
is combinatorially rigid if its combinatorial class consists of a single
quadratic polynomial.
\end{definition}

\begin{definition}[Removable compact set; robust polynomial]
\label{def:removable-robust}
A compact set $K\subset\overline{\C}$ is removable if every conformal embedding of
$U\setminus K$ extends conformally across $K$.  A map is robust if, for arbitrarily
large $m$, there is an annulus in $\C\setminus\omega(c)$ of definite modulus,
homotopic rel.\ $\omega(c)$ to a Jordan curve enclosing $J(R^mf)$ and containing
no point of $\omega(c)\setminus J(R^mf)$.  Here $\omega(c)$ is the $\omega$-limit
set of the critical point.
\end{definition}

\subsection*{Results}

\begin{lemma}[Lemma~0]
\label{lem:lyubich-zero}
$V^1\Subset V^0$.  Status: proved in~\cite{LyubichGeometry}.
\end{lemma}

\begin{theorem}[Lyubich, Theorem~I]
\label{thm:lyubich-I}
Let $f$ be quadratic-like with definite modulus and with internal class in a
fixed finite family of truncated secondary limbs.  Let $n(k)$ count the levels
with non-central return that precede the $k$-th quadratic-like level.  Then
$\mu_{n(k)+1}=\modulus(V^{n(k)}\setminus V^{n(k)+1})$ grows uniformly linearly
in $k$.  Status: proved in the source (Section~3), using Lemmas~1--10 and cited
distortion and geometry estimates.
\end{theorem}

\begin{theorem}[Lyubich, Theorem~II]
\label{thm:lyubich-II}
For every $f\in\mathcal F$, all quadratic-like renormalizations $R^nf$ have
definite moduli.  Status: proved in the source from Theorem~\ref{thm:lyubich-I}
and Lemma~11.
\end{theorem}

\begin{corollary}[Lyubich, Corollary~IIa]
\label{cor:lyubich-IIa}
For $f\in\mathcal F$ the set $\omega(c)$ is a removable Cantor set.  Status:
derived in the source from Theorem~\ref{thm:lyubich-II} and an annulus-removability
criterion.
\end{corollary}

\begin{corollary}[Lyubich, Corollary~IIb]
\label{cor:lyubich-IIb}
For $f\in\mathcal F$ the Julia set $J(f)$ is locally connected.  Status: derived
in the source from Theorem~\ref{thm:lyubich-II}, citing Hu--Jiang and Jiang.
\end{corollary}

\begin{corollary}[Lyubich, Corollary~IIc]
\label{cor:lyubich-IIc}
Every $f\in\mathcal F$ is robust.  Status: derived in the source from
Theorem~\ref{thm:lyubich-II}.  The implication ``robust $\Rightarrow$ no invariant
line fields'' is cited to McMullen.
\end{corollary}

\begin{theorem}[Lyubich, Theorem~III]
\label{thm:lyubich-III}
If $f,\tilde f\in\mathcal F$ are combinatorially equivalent, then they are
quasiconformally conjugate.  Status: proved in the source by the pull-back
argument of its Section~4.
\end{theorem}

\begin{corollary}[Lyubich, combinatorial rigidity]
\label{cor:lyubich-rigidity}
Every quadratic polynomial $z^2+c$ in $\mathcal F$ is combinatorially rigid.
Consequently MLC holds at $c$.  Status: proved in the source from
Theorem~\ref{thm:lyubich-III} and the standard argument that a closed
combinatorial class is either open or a singleton quasiconformal class.
\end{corollary}

\begin{theorem}[Lyubich, Theorem~D; real maps]
\label{thm:lyubich-D}
For the class of real maps of fixed $q$-renormalizable type with initial modulus
at least $\bar\mu$, either the renormalized modulus satisfies
$\modulus(Rf)\ge\bar\mu$, or the real geometry is essentially bounded until the
next quadratic-like level.  Status: proved in Section~5 with Sullivan's and real
bounds as inputs.
\end{theorem}

\subsection*{Dependencies and connections}

The source uses:
the Douady--Hubbard straightening theorem, to identify each hybrid class with a unique
polynomial representative~\cite{DHPolynomialLike};
Yoccoz's puzzle and shrinking theorem, for non-renormalizable nests~\cite{HubbardYoccoz};
Sullivan's complex bounds and sector lemma;
Hu--Jiang and Jiang, for local connectivity from a priori bounds;
McMullen, for robustness and the absence of invariant line fields;
and Lyubich's own papers on generalized polynomial-like maps, real geometry, and
non-quadratic-like complex bounds~\cite{LyubichDynamics}.

Theorem~\ref{thm:lyubich-I} is the analytic core: definite moduli are preserved
under renormalization in the class.  Theorem~\ref{thm:lyubich-III} transfers
combinatorial information to conformal information, and the rigidity corollary
gives MLC at the corresponding parameters.  The later sources replace the
hypotheses (C1), (C2) by other combinatorial conditions; see
Chapter~\ref{ch:digests-modern}.

\subsection*{Formal status}

No statement of this section is formalized.  The repository has no conformal
modulus, no quasiconformal conjugacy, and no straightening theorem.  Lyubich's
hypotheses appear in the repository only as interface data, as described in
Chapter~\ref{ch:conditional-matrix} and Appendix~\ref{ch:formal-map}.
